\chapter{Modelos de colas}

\section{Proposito}
La teoria de colas estudia sistemas donde clientes, trabajos o unidades esperan
servicio ante capacidad limitada. Sus metricas centrales incluyen utilizacion,
longitud esperada de cola, tiempo esperado en el sistema y probabilidad de
bloqueo \cite{gross2008fundamentals}.

\section{Modelo M/M/1}
Para llegadas Poisson con tasa \(\lambda\), tiempos de servicio exponenciales
con tasa \(\mu\) y un servidor, si \(\rho=\lambda/\mu<1\), entonces
\[
  L = \frac{\rho}{1-\rho}, \qquad
  L_q = \frac{\rho^2}{1-\rho}, \qquad
  W = \frac{1}{\mu-\lambda}, \qquad
  W_q = \frac{\lambda}{\mu(\mu-\lambda)}.
\]

\section{Criterio de documentacion}
Los ejemplos de colas en \qsb{} deberan registrar la notacion Kendall, la unidad
de tiempo, las distribuciones asumidas, la tasa efectiva de llegada y las notas
del backend cuando se use una aproximacion.
